% Préambule LaTeX - Configuration du document
% Conforme au template officiel UNamur Master Sciences Informatiques

% ============================================================================
% PACKAGES DE BASE
% ============================================================================

% Encodage et langue
\usepackage[utf8]{inputenc}
\usepackage[T1]{fontenc}
\usepackage[sfdefault]{atkinson}  % Police Atkinson Hyperlegible (obligatoire UNamur)
\usepackage[french]{babel}

% Marges et mise en page (conformes template UNamur)
\usepackage[top=2cm, bottom=2.5cm, left=2cm, right=2cm]{geometry}
\usepackage{setspace}
\onehalfspacing  % Interligne 1.5

% Polices et typographie
\renewcommand{\familydefault}{\sfdefault}
\usepackage{microtype}  % Amélioration typographie

% ============================================================================
% PACKAGES GRAPHIQUES
% ============================================================================

\usepackage{graphicx}
\usepackage{float}
\usepackage[dvipsnames]{xcolor}
\usepackage{tikz}  % Pour page de couverture
\usepackage{fancybox}

% Chemin vers les figures
\graphicspath{{figures/}{img/}}

% ============================================================================
% PACKAGES MATHÉMATIQUES
% ============================================================================

\usepackage{amsmath}
\usepackage{amssymb}
\usepackage{amsthm}

% ============================================================================
% PACKAGES TABLEAUX
% ============================================================================

\usepackage{array}
\usepackage{booktabs}
\usepackage{longtable}
\usepackage{multirow}

% ============================================================================
% PACKAGES LISTES ET ÉNUMÉRATIONS
% ============================================================================

\usepackage{enumitem}
\usepackage{pifont}

% ============================================================================
% PACKAGES BIBLIOGRAPHIE
% ============================================================================

\usepackage[backend=biber,style=numeric,sorting=none,maxbibnames=99]{biblatex}
\addbibresource{bibliography.bib}

% ============================================================================
% PACKAGES LIENS ET RÉFÉRENCES
% ============================================================================

\usepackage[hidelinks]{hyperref}
\usepackage{url}
\usepackage{cleveref}

\hypersetup{
    colorlinks=false,
    linkcolor=black,
    citecolor=black,
    urlcolor=blue,
    pdftitle={Mesures DNS dans l'espace et le temps},
    pdfauthor={Olivier Gautier},
    pdfsubject={Master 60 Sciences Informatiques - UNamur},
    pdfkeywords={DNS, RIPE Atlas, mesures distribuées, géolocalisation}
}

% Configuration des URLs
\urlstyle{same}

% ============================================================================
% PACKAGES CODE SOURCE
% ============================================================================

\usepackage{listings}
\usepackage[linesnumbered,ruled,vlined]{algorithm2e}

% Configuration listings - Couleurs
\definecolor{codebg}{rgb}{0.97,0.97,0.97}
\definecolor{javakeyword}{RGB}{0,0,180}
\definecolor{javacomment}{RGB}{0,128,0}
\definecolor{javastring}{RGB}{163,21,21}

% Style Python
\lstdefinestyle{python}{
  language=Python,
  backgroundcolor=\color{codebg},
  basicstyle=\ttfamily\small,
  keywordstyle=\color{blue}\bfseries,
  commentstyle=\color{gray}\itshape,
  stringstyle=\color{orange},
  showstringspaces=false,
  numbers=left,
  numberstyle=\tiny\color{gray},
  numbersep=10pt,
  frame=single,
  tabsize=4,
  captionpos=b,
  breaklines=true,
  breakatwhitespace=true
}

% Style Bash
\lstdefinestyle{bash}{
  language=bash,
  backgroundcolor=\color{codebg},
  basicstyle=\ttfamily\small,
  keywordstyle=\color{blue}\bfseries,
  commentstyle=\color{gray}\itshape,
  stringstyle=\color{orange},
  showstringspaces=false,
  numbers=left,
  numberstyle=\tiny\color{gray},
  numbersep=10pt,
  frame=single,
  breaklines=true
}

% Style JSON
\lstdefinelanguage{json}{
  morestring=[b]",
  morestring=[d]',
  literate=
   *{0}{{{\color{blue}0}}}{1}
    {1}{{{\color{blue}1}}}{1}
    {2}{{{\color{blue}2}}}{1}
    {3}{{{\color{blue}3}}}{1}
    {4}{{{\color{blue}4}}}{1}
    {5}{{{\color{blue}5}}}{1}
    {6}{{{\color{blue}6}}}{1}
    {7}{{{\color{blue}7}}}{1}
    {8}{{{\color{blue}8}}}{1}
    {9}{{{\color{blue}9}}}{1}
    {:}{{{\color{red}:}}}{1}
    {,}{{{\color{red},}}}{1}
    {\{}{{{\color{red}\{}}}{1}
    {\}}{{{\color{red}\}}}}{1}
    {[}{{{\color{red}[}}}{1}
    {]}{{{\color{red}]}}}{1}
    {true}{{{\color{purple}true}}}{4}
    {false}{{{\color{purple}false}}}{5}
    {null}{{{\color{purple}null}}}{4},
}

\lstdefinestyle{json}{
  language=json,
  backgroundcolor=\color{codebg},
  basicstyle=\ttfamily\small,
  keywordstyle=\color{blue},
  stringstyle=\color{orange},
  showstringspaces=false,
  numbers=left,
  numberstyle=\tiny\color{gray},
  numbersep=10pt,
  frame=single,
  breaklines=true,
  captionpos=b
}

% ============================================================================
% EN-TÊTES ET PIEDS DE PAGE
% ============================================================================

\usepackage{fancyhdr}

\pagestyle{fancy}
\fancyhf{}
\fancyhead[L]{\nouppercase{\leftmark}}
\fancyhead[R]{\thepage}
\renewcommand{\headrulewidth}{1pt}
\renewcommand{\footrulewidth}{0pt}

% Style pour pages de chapitre
\fancypagestyle{plain}{
    \fancyhf{}
    \fancyfoot[C]{\thepage}
    \renewcommand{\headrulewidth}{0pt}
}

% Hauteur en-tête
\setlength{\headheight}{13.59999pt}

% ============================================================================
% PACKAGES CAPTIONS
% ============================================================================

\usepackage{caption}
\usepackage{subcaption}

\captionsetup{
    font=small,
    labelfont=bf,
    textfont=it
}

% ============================================================================
% CONFIGURATION DOCUMENT
% ============================================================================

% Profondeur table des matières
\setcounter{tocdepth}{2}
\setcounter{secnumdepth}{3}

% Espacement paragraphes
\setlength{\parskip}{0.5em}
\setlength{\parindent}{0pt}

% Justification
\usepackage{ragged2e}

% ============================================================================
% COMMANDES PERSONNALISÉES
% ============================================================================

% Commande pour le symbole de référence
\newcommand{\refmark}[1]{\textsuperscript{\cite{#1}}}

% Commande pour faciliter les références
\newcommand{\figref}[1]{Figure~\ref{#1}}
\newcommand{\tabref}[1]{Tableau~\ref{#1}}
\newcommand{\chapref}[1]{Chapitre~\ref{#1}}
\newcommand{\secref}[1]{Section~\ref{#1}}

% ============================================================================
% FIN DU PRÉAMBULE
% ============================================================================
